\documentclass[12pt,a4paper]{article}
\usepackage[a4paper, margin=2.5cm]{geometry}
\usepackage[utf8]{inputenc}
\usepackage[T1]{fontenc}
\usepackage[magyar]{babel}
\usepackage{amsmath, amssymb, amsthm}
\usepackage{xcolor}
\usepackage{tcolorbox}
\usepackage{enumitem}

\title{Ismétlés nélküli permutáció - Megoldókulcs}
\author{Részletes levezetések és érettségi pontozási útmutató}
\date{}

%%% Megjelenítéshez automatikusan kiegészített csomagok %%%
\usepackage{tikz}
\usepackage{pgfplots}
\pgfplotsset{compat=1.18}
\usepackage{graphicx}
\usepackage{booktabs}
\usepackage{float}
\usepackage{caption}
\usepackage{microtype}

\usepackage{hyperref}

\begin{document}

\maketitle

\section*{Részletes megoldások és pontozási útmutató}

\begin{enumerate}[label=\arabic*., leftmargin=*]
    \item \textbf{1. Feladat megoldása}
    \begin{enumerate}[label=\alph*)]
        \item \textbf{Kifejezések kiszámítása:}
        \begin{align*}
            \frac{6!}{4!} &= \frac{6 \cdot 5 \cdot 4!}{4!} = 6 \cdot 5 = 30 \\
            \frac{7!}{5! \cdot 2!} &= \frac{7 \cdot 6 \cdot 5!}{5! \cdot (2 \cdot 1)} = \frac{42}{2} = 21
        \end{align*}
        \item \textbf{Könyvek sorba rendezése:}
        Mivel 4 különböző könyvről van szó, és mind a 4-et sorba rakjuk:
        \begin{equation*}
            P_4 = 4! = 4 \cdot 3 \cdot 2 \cdot 1 = 24
        \end{equation*}
        \textbf{Válasz:} 24-féleképpen rendezhetjük sorba a regényeket.
    \end{enumerate}
    \textit{Pontozás:} a) alpont 2 pont (egyenként 1-1); b) alpont 2 pont (képlet felírása 1 pont, végeredmény 1 pont). Összesen: 4 pont.

    \vspace{0.4cm}

    \item \textbf{2. Feladat megoldása}
    \begin{enumerate}[label=\alph*)]
        \item \textbf{Összes lehetséges végeredmény:} 
        5 különböző futó sorrendje:
        \begin{equation*}
            P_5 = 5! = 120
        \end{equation*}
        \item \textbf{Rögzített első két hely:}
        Ha Anna az 1. és Bence a 2., akkor a maradék 3 helyre a maradék 3 futó (Cecília, Dániel, Eszter) tetszőleges sorrendben ülhet be:
        \begin{equation*}
            P_3 = 3! = 3 \cdot 2 \cdot 1 = 6
        \end{equation*}
    \end{enumerate}
    \textbf{Válaszok:} a) 120 féle; b) 6 féle végeredmény létezik.
    \textit{Pontozás:} a) alpont 2 pont; b) alpont 3 pont (feltétel felismerése 2 pont, számolás 1 pont). Összesen: 5 pont.

    \vspace{0.4cm}

    \item \textbf{3. Feladat megoldása}
    \begin{enumerate}[label=\alph*)]
        \item \textbf{Korlátozás nélküli eset:} 6 ember a 6 széken: $6! = 720$.
        \item \textbf{Lányok egy blokkban (ragasztós trükk):} 
        A 3 lányt egyetlen egységnek (csomagnak) tekintjük. Így a fiúkkal együtt 4 egységünk van, ezek sorrendje $4! = 24$. A lányok a saját blokkjukon belül $3! = 6$-féleképpen cserélhetnek helyet.
        \begin{equation*}
            4! \cdot 3! = 24 \cdot 6 = 144
        \end{equation*}
        \item \textbf{Fiúk a két szélső széken:} 
        A 3 fiú közül 2-t kiválasztunk a két szélső székre, és sorba rendezzük őket: $V_3^2 = 3 \cdot 2 = 6$ lehetőség. A maradék 4 ember a belső 4 székre tetszőlegesen leülhet: $4! = 24$ lehetőség. Összesen: $6 \cdot 24 = 144$.
    \end{enumerate}
    \textit{Pontozás:} a) alpont 1 pont; b) alpont 3 pont (blokkolás 2 pont, belső helycsere 1 pont); c) alpont 3 pont (szélső helyek kezelése 2 pont, belső permutáció 1 pont). Összesen: 7 pont.

    \vspace{0.4cm}

    \item \textbf{4. Feladat megoldása}
    \begin{enumerate}[label=\alph*)]
        \item \textbf{5-jegyű számok készítése:} 5 különböző számjegy ($1, 3, 5, 7, 9$):
        \begin{equation*}
            P_5 = 5! = 120
        \end{equation*}
        \item \textbf{Oszthatóság 5-tel:} 
        Egy szám akkor osztható 5-tel, ha az utolsó számjegye 0 vagy 5. Mivel csak az 5-ös szerepel a rendelkezésre álló számjegyek között (0 nincs), ezért \textbf{minden egyes képzett szám} az 5-re végződik! 
        Tehát a feltételnek mind a 120 szám eleget tesz.
    \end{enumerate}
    \textit{Pontozás:} a) alpont 2 pont; b) alpont 3 pont (oszthatósági szabály ismerete 1 pont, felismerés, hogy az összes szám ilyen 2 pont). Összesen: 5 pont.

    \vspace{0.4cm}

    \item \textbf{5. Feladat megoldása}
    \begin{enumerate}[label=\alph*)]
        \item \textbf{Elnök és titkár nem ülhet egymás mellett (komplementer módszer):}
        \begin{itemize}
            \item Összes eset: $8! = 40320$.
            \item Tiltott esetek (egymás mellett vannak): Az elnököt és titkárt összeragasztjuk (2! helycsere), így 7 egységünk van ($7!$). Tiltott esetek száma: $7! \cdot 2! = 5040 \cdot 2 = 10080$.
            \item Megengedett esetek: $40320 - 10080 = 30240$.
        \end{itemize}
        \item \textbf{Pontosan 2 személy ül közöttük:}
        \begin{itemize}
            \item Kiválasztjuk a 2 közbülső személyt a maradék 6-ból: $V_6^2 = 6 \cdot 5 = 30$ vagy kombinációval $ \binom{6}{2} \cdot 2! = 30 $.
            \item Ezt a 4 fős blokkot (elnök + 2 ember + titkár) egy egységként kezeljük, melyen belül az elnök és a titkár helyet cserélhet (2 lehetőség).
            \item A többi 4 emberrel együtt ez 5 egység, melyek sorrendje $5! = 120$.
            \item Összesen: $30 \cdot 2 \cdot 120 = 7200$.
        \end{itemize}
    \end{enumerate}
    \textit{Pontozás:} a) alpont 4 pont (összes eset 1p, tiltott esetek 2p, kivonás 1p); b) alpont 5 pont. Összesen: 9 pont.

    \vspace{0.4cm}

    \item \textbf{6. Feladat megoldása (Emelt szintű kihívás)}
    A feladat két korlátozást tartalmaz: a hidrogén (H) és oxigén (O) nem lehet egymás mellett, valamint a vas (Fe) és réz (Cu) sem állhat közvetlenül egymás mellett.
    \begin{itemize}
        \item \textbf{Összes elhelyezés:} 7 különböző elem $\rightarrow$ $7! = 5040$.
        \item \textbf{H és O együtt van:} Blokkolva $6! \cdot 2! = 720 \cdot 2 = 1440$ eset.
        \item \textbf{Fe és Cu együtt van:} Blokkolva $6! \cdot 2! = 1440$ eset.
        \item \textbf{Mindkét páros egyszerre együtt van:} Ha H-O egy blokk ($\text{BO}$) és Fe-Cu is egy blokk ($\text{BC}$), valamint a maradék 3 elem ($M_1, M_2, M_3$), akkor összesen 5 egységünk van ($5! = 120$). Mivel mindkét blokkon belül lehet helycsere ($2! \cdot 2! = 4$), a metszet mérete: $5! \cdot 2! \cdot 2! = 120 \cdot 4 = 480$.
        \item \textbf{Inklúzió-exklúzió elve alapján a tiltott esetek száma:}
        \begin{equation*}
            \text{Tiltott} = 1440 + 1440 - 480 = 2400
        \end{equation*}
        \item \textbf{Végső megengedett darabszám:}
        \begin{equation*}
            \text{Megengedett} = 5040 - 2400 = 2640
        \end{equation*}
    \end{itemize}
    \textbf{Válasz:} A tanár 2640-féleképpen helyezheti el a modelleket.
    \textit{Pontozás:} Összes eset meghatározása (1p), külön-külön tiltott esetek kiszámítása (2p), metszet kiszámítása (2p), inklúzió-exklúzió helyes alkalmazása (2p), végeredmény (1p). Összesen: 8 pont.
\end{enumerate}

\end{document}
EOF